% ===== Mục 4.4 =====
\subsubsection{Nắn hình học và đo chất lượng ảnh}
Bảng~\ref{tab:D2} so sánh các cách tìm 4 góc trong thí nghiệm D2.
Một mẫu được tính là thành công khi sai số trung bình của 4 góc nhỏ hơn 2\% độ dài đường chéo ảnh.
Ngưỡng Canny chọn trên tập kiểm định là $(T_2,T_1)=(30,60)$, đúng tỉ lệ $T_1=2T_2$ khuyến nghị trong học phần \chuong{5, S47}.

\begin{table}[H]
\centering
\caption{Tìm 4 góc trên 210 mẫu có phối cảnh của tập kiểm tra \bo{5} (thí nghiệm D2).}\label{tab:D2}
\small
\begin{tabular}{lrrr}
\toprule
\textbf{Phương án} & \textbf{Tìm được (\%)} & \textbf{Thành công (\%)} & \textbf{Sai số trung vị (px)} \\
\midrule
Canny, đường bao, xấp xỉ tứ giác & 99,0 & \textbf{99,0} & \textbf{1,43} \\
Hough & 87,1 & 84,8 & 3,63 \\
Hình chữ nhật bao nhỏ nhất & \textbf{100} & 59,5 & 26,4 \\
Kết hợp có dự phòng (pipeline) & \textbf{100} & \textbf{99,0} & \textbf{1,43} \\
\bottomrule
\end{tabular}
\end{table}

Tìm góc bằng Canny và đường bao thành công ở 99,0\% số mẫu với sai số trung vị 1,43~px.
Hình chữ nhật bao nhỏ nhất luôn trả về một khung, nhưng chỉ thành công ở 59,5\% số mẫu.
Nguyên nhân là hình chữ nhật không mô tả được tứ giác bị méo phối cảnh.

\begin{figure}[H]
\centering
\includegraphics[width=0.5\linewidth]{preprocessing/D3_curve.png}
\caption{Sai số tìm 4 góc theo góc nghiêng và mức phối cảnh trên tập kiểm tra \bo{5} (thí nghiệm D3, lưới 15 ô, mỗi ô 10 ảnh).}\label{fig:D3}
\end{figure}

Hình~\ref{fig:D3} cho thấy sai số tìm góc tăng theo góc nghiêng.
Sai số lớn nhất là 2,3~px, ở góc nghiêng 30° kết hợp phối cảnh 20\%, với 8/10 ảnh thành công.
Nguyên nhân là ở mức này một góc của tờ giấy nằm sát mép ảnh.

Trong thí nghiệm D4, chỉnh nghiêng bằng lược đồ chiếu~\eqref{eq:skew} có sai số $0^\circ$ trên 120 ảnh chỉ bị xoay (từ $\pm2^\circ$ đến $\pm15^\circ$), với thời gian 79~ms.
Biến đổi Radon cũng có sai số $0^\circ$ nhưng chậm hơn 120 lần (9,6~s).
Biến đổi Hough có sai số trung bình $0{,}06^\circ$ (Bảng~\ref{tab:D4}, Phụ lục~\ref{app:preprocessing}).
Trong thí nghiệm D1, độ đo mờ $\sigma^2_L$ theo~\eqref{eq:blur} đạt AUC \num{0.979}, cao hơn Tenengrad (\num{0.954}) và năng lượng tần số cao (\num{0.920}).

Ca lỗi điển hình là ảnh mẫu \texttt{rx\_bo5\_degraded\_01}, bị nghiêng, méo phối cảnh, có bóng và bị mờ.
Ảnh này được nắn thẳng hoàn toàn, nhưng độ mờ làm bước (4) không còn phân biệt được dòng đậm (Mục~\ref{sec:e-pres}).

\subsubsection{Phân loại tài liệu}
Bảng~\ref{tab:D10} so sánh ba bộ đặc trưng cho bộ phân loại kNN có từ chối trong thí nghiệm D10.
Hệ số $k$ và độ đo khoảng cách được chọn bằng kiểm định chéo 5-fold trên 439 ảnh huấn luyện.

\begin{table}[H]
\centering
\caption{Phân loại 3 lớp bằng kNN có từ chối (thí nghiệm D10).
Tập kiểm tra cùng nguồn gồm 218 ảnh của \bo{1}, \bo{2} và \bo{3}.
Tập mở gồm 100 ảnh thuộc các lớp RVL-CDIP không có khi huấn luyện.
\Bo{4} gồm 83 ảnh chụp thật.}\label{tab:D10}
\small
\begin{tabular}{lccc}
\toprule
\textbf{Đặc trưng} & \makecell{\textbf{Cùng nguồn:}\\\textbf{Acc / Macro-F1}} & \makecell{\textbf{Tập mở:}\\\textbf{bị từ chối (\%)}} & \makecell{\textbf{\Bo{4}:}\\\textbf{Acc / Macro-F1}} \\
\midrule
Histogram màu $8^3$ bin (đối chứng) & 0,927 / 0,801 & 95 & \textbf{0,494 / 0,349} \\
HOG, GLCM và bố cục trên ảnh CLAHE & \textbf{0,963 / 0,930} & 90 & 0,301 / 0,251 \\
HOG và bố cục trên ảnh chuẩn hóa (pipeline) & 0,936 / 0,854 & \textbf{97} & 0,386 / 0,292 \\
\bottomrule
\end{tabular}
\end{table}

Trên dữ liệu cùng nguồn, cả ba bộ đặc trưng đều đạt độ chính xác trên 92\% và từ chối được 90--97\% tài liệu lạ.
Trên ảnh chụp thật của \bo{4}, độ chính xác của cả ba bộ đặc trưng giảm xuống còn 0,30--0,49.
Nguyên nhân là mô hình học đặc điểm ảnh của từng nguồn thay vì bố cục tài liệu.
Ảnh VAIPE là ảnh scan, ảnh \bo{2} bị ép về $640\times640$, còn ảnh RVL-CDIP là ảnh scan đen trắng.
Chuẩn hóa ảnh bằng khử bóng và nhị phân hóa nâng độ chính xác của HOG từ 0,30 lên 0,39.
Ma trận nhầm lẫn cho thấy cả 18 đơn thuốc và 33/48 phiếu xét nghiệm thật bị từ chối thành lớp ``khác'', vì chúng khác xa mọi mẫu huấn luyện.
Đây là hạn chế chính của \bo{1}--\bo{3}.
Hướng khắc phục là bổ sung ảnh chụp thật vào tập huấn luyện, không lấy từ \bo{4} vì \bo{4} là tập kiểm tra.
Trong bản trình diễn, người dùng tự chọn loại tài liệu khi bộ phân loại từ chối.
